\section{Axiomatised interfaces: how FV-Spec absorbs effectful source PBTs}

A reviewer will reasonably ask: \emph{the source PBTs come from real
production code. Real production code touches Redis, sockets, the
filesystem, the wall clock, GitHub webhooks. How can the resulting Lean
theorems possibly say anything meaningful about that?} The answer is
that the formalisation agent treats the effectful world as an
\emph{uninterpreted interface}: the parts of the source that are
essential to the property under test are translated into Lean
structures and inductive types, and the parts that are mere plumbing
--- the database, the queue, the system clock --- are introduced as
opaque \lstinline{axiom} declarations and threaded through the theorem
statements as parameters. A theorem holds \emph{over all
interpretations} of those axioms, which is exactly the contract the
original PBT was checking against the live environment.

A canonical instance is \texttt{test\_workflow\_job\_time\_to\_start},
drawn from \texttt{scality/runner-manager}\footnote{\url{https://github.com/scality/runner-manager}}
(a GitHub Actions self-hosted runner controller). The Python PBT
initialises Redis-backed model classes, runs an ORM migrator, calls
\lstinline{datetime.now()} to fabricate two timestamps straddling a
runner-startup timeout, enqueues a webhook through a background job
queue, and asserts on the resulting runner count. The fvspec
translation introduces seventeen axioms --- including
\lstinline{axiom Redis : Type}, \lstinline{axiom Queue : Type},
\lstinline{axiom State : Type},
\lstinline{axiom State.after_enqueue : State → RunnerGroup → WorkflowJobEvent → ExtendedSettings → State},
and
\lstinline{axiom get_runners : RunnerGroup → State → List Runner}
--- and then states the behavioural contract purely in terms of these
uninterpreted state-transition operators:

\begin{lstlisting}[language=Lean,basicstyle=\ttfamily\small,frame=single]
theorem create_runner_when_above_timeout
  (...)
  (h_initial_count : (get_runners runner_group initial_state).length = 1)
  (h_below_max    : 1 < runner_group.max)
  (h_above_timeout : started_at - created_at > settings.timeout_runner)
  (h_enqueue : state_after_enqueue =
    State.after_enqueue initial_state runner_group webhook_updated settings)
  : (get_runners runner_group state_after_enqueue).length = 2 := by
  sorry
\end{lstlisting}

This is genuinely different from what a benchmark drawn from Mathlib,
HumanEval, or a pedagogical FV corpus can ask of a model. Those
corpora pre-select for problems that are \emph{self-contained by
construction}: a number-theoretic identity, a small algorithm with no
side-effects, a closed-world data structure. FV-Spec's source
distribution does the opposite --- it pre-selects for properties that
practitioners actually wanted to test, which means the theorems carry
the shape of real software: state machines, external services, time,
concurrency. The axiomatisation step is what converts that shape into
a well-formed proof obligation without losing it. A model proving
\texttt{create\_runner\_when\_above\_timeout} is reasoning about a
scheduler's state-transition semantics under a timeout precondition;
the fact that the underlying scheduler in production talks to Redis is
no more load-bearing in the theorem than the choice of register
allocator is in a correctness proof of a sorting algorithm. We claim
this is the honest way to ship FV evaluation problems sourced from real
code, and that no existing FV benchmark does it because none of them
source from real code in the first place.
